% !TeX root = ../教材题库_全部条目_转写工作稿.tex
% Source: XJ-B1-C05-EX-P190-Y1515
% 第5章 / 5.4 函数 y=A sin(ωx+φ) 的图象与性质 / 教材 p.190、191 / PDF 第197、198页
% OCR未校对：公式与文字均可能有误；未完成数学TeX转写；禁止直接发给学生
% 原题及相邻原解截图（核对用）：../题目原图/XJ-B1-C05-EX-P190-Y1515-01.webp, ../题目原图/XJ-B1-C05-EX-P190-Y1515-02.webp
% 本文件只在转写完成后作为题干使用；OCR 条目尚不是可靠题干。
\begin{bankitem}{5.4 例题 例题3}{第5章 / 5.4 函数 y=A sin(\ensuremath{\omega}x+\ensuremath{\varphi}) 的图象与性质 / 教材 p.190、191 / PDF 第197、198页}{XJ-B1-C05-EX-P190-Y1515}{OCR未校对：公式与文字均可能有误；未完成数学TeX转写；禁止直接发给学生}
\begin{ocrdraft}
[教材p.190 / PDF第197页]
\par
0河作出函数y一3sin—x
\par
在长度为一个周期的闭区间上的简图，并说明3-
\par
的图象是由函数的图象经过怎样的变化而得到的\ensuremath{\cdot}
\par
3Sin\ocrglyph{00AF}X
\par
[教材p.191 / PDF第198页]
\par
2丌
\par
4丌，0，4
\par
“〕上求出函数3，
\par
解函数3，=3sin一工的周期为
\par
=3sin
\par
—x图
\par
象上的五个关键点，列表如下：
\par
3丌
\par
），一3sin一工
\par
“〕的大致图象，
\par
描点，连线，作出函数
\par
如图5．4一6．
\par
工，xe0，
\par
y=3s1n\ocrglyph{00E4}X
\par
y=sln\ocrglyph{00E4}X
\par
4兀
\par
y=sln
\par
0
\par
丌2丌
\par
2丌X
\par
图5．4一6
\par
观察图5．4一6可以发现，函数y=3sin一的图象可以由函数，一x的图象
\par
经过下述变化得到：先把函数3，一sin的图象上每一点的纵坐标不变、横坐标缩短
\par
到原来的，再把所得图象上各点的横坐标不变、纵坐标伸长到原来的3倍，即得
\par
到函数3，=3sin—x的图象．
\end{ocrdraft}
\end{bankitem}
